\scanpage{53}
Pour \uline{l'existence}, considérons l'homomorphisme canonique
\[
\varphi_X:X\longrightarrow A_X^1=X'
\]
et le cycle
\[
Z'=(\varphi_X\times\mathrm{id}_Y)_*(Z)\in\mathrm{Cyc}_m(X'\times Y)
\]
Introduisons également
\[
B=B_X=\underline{\mathrm{Pic}}^{00}_{A_X^1/k}\simeq\underline{\mathrm{Pic}}^0_{A_X/k}
\]
(isomorphe canoniquement à $\underline{\mathrm{Pic}}^{00}_{X/k}$ par l'hom $\underline{\mathrm{Pic}}^0_{A_X^1/k}\unsure{-|\longrightarrow}\underline{\mathrm{Pic}}^0_{X/k}$ transposé de $\varphi_X$). Considérons le cycle $Z''=B\times Z'$ sur $B\times X'\times Y$, et l'élément canonique $\Theta$ de $\mathrm{Pic}(B\times X')$, dont l'image inverse \deleted{\unsure{dans}} $B\times X'\times Y$ \unsure{sera|ser.} notée $\Theta'$. Soit $\deleted{\unsure{\theta|\Theta'}}\theta'$ la classe de diviseurs, mod équivalence numérique, définie par $\Theta'$, et soit de $\overline{\mathrm{m}}$ $z''$ la classe d'équivalence numérique de $Z''$ (qui ne dépend, bien entendu, que de la classe d'équiv. numérique de $Z$). Soit \deleted{$\unsure{\pi:}$}
\[
\pi:B\times X'\times Y\longrightarrow B\times Y
\]
la projection, et formons
\[
c=\pi_*(z''\cdot\unsure{\Theta|\theta})
\]
qui est une classe d'équiv \uline{numérique} de diviseurs sur $B\times Y$, \uline{donc définit} une classe
% circled note
\margin{\unsure{une référence|en référence}}
% archive sheet number: 52
